% !TeX root = ../Book4.tex
\begin{figure}[htbp]
	\centering
	\begin{tikzpicture}[x=1.5cm,y=1.2cm,>=Stealth,font=\small]
		\draw[->] (-0.45,0)--(2.7,0) node[right] {\(q\)};
		\draw[->] (0,-1.6)--(0,1.65) node[above] {\(p\)};
		\draw[GeometryBlue,thick,domain=-1.35:1.35,samples=80,smooth,variable=\t] plot ({\t*\t},\t);
		\draw[dashed,GeometryGray] (1,-1)--(1,1);
		\draw[->,GeometryBlue] (1,0.84)--(1,0.12);
		\draw[->,GeometryBlue] (1,-0.84)--(1,-0.12);
		\fill[GeometryBlue] (1,1) circle (2pt) node[above left] {\((1,1)\)};
		\fill[GeometryBlue] (1,-1) circle (2pt) node[below left] {\((1,-1)\)};
		\fill (1,0) circle (1.6pt) node[below right] {\(1\)};
		\fill (0,0) circle (1.6pt) node[below left] {\(0\)};
		\node[right] at (1.9,1.25) {\(q=p^2\)};
	\end{tikzpicture}
	\caption[Lagrange 曲线与底投影]{光滑曲线 \(q=p^2\) 向 \(q\) 轴投影时，两点可落到同一个底点，原点处的切线还完全竖直。曲线的光滑性并不保证它是关于 \(q\) 的单值图像。}
	\label{g4:fig:0804:graph-projection}
\end{figure}
